\begin{examplebox}
\textbf{例题1：} 写出命题"若 \(x>0\)，则 \(x^2>0\)"的逆命题、否命题、逆否命题，并判断它们的真假。

\textit{解析：}
\begin{itemize}
    \item 原命题：真（正数的平方为正）。
    \item 逆命题：若 \(x^2>0\)，则 \(x>0\)。假（如 \(x=-1\) 时 \(x^2>0\) 但 \(x\not>0\)）。
    \item 否命题：若 \(x\le 0\)，则 \(x^2\le 0\)。假（如 \(x=-1\) 时 \(x\le0\) 但 \(x^2=1>0\)）。
    \item 逆否命题：若 \(x^2\le0\)，则 \(x\le0\)。真（只有 \(x=0\) 时 \(x^2\le0\) 成立，此时 \(x\le0\) 成立）。
\end{itemize}
注意原命题与逆否命题同真，逆命题与否命题同假。
\end{examplebox}

\begin{examplebox}
\textbf{例题2：} 已知命题"若 \(m>0\)，则方程 \(x^2+x-m=0\) 有实根"为真，写出它的逆否命题并判断真假。

\textit{解析：} 逆否命题：若方程 \(x^2+x-m=0\) 无实根，则 \(m\le 0\)。由原命题为真，逆否命题必为真（等价性）。可直接验证：方程无实根 \(\Leftrightarrow \Delta=1+4m<0\Leftrightarrow m<-\frac14\)，此时 \(m\le0\) 成立，故逆否命题为真。
\end{examplebox}